\section{Related work}
\label{app:related}

\paragraph{Instruction-following evaluation.}
IFEval, IFBench, InFoBench, and ComplexBench \citep{zhou_instruction-following_2023, pyatkin_generalizing_2025, qin_infobench_2024, wen_benchmarking_2024} and similar benchmarks measure whether models can follow   explicit constraints of varying complexity, but none introduce conflicting in-context evidence that pressures the model to abandon its instruction. Our paradigm complements these benchmarks by isolating the induction-instruction axis. Compared to existing work, our setup differs in three key respects: the instructions are intentionally simple (removing instruction complexity as a confound), the challenge arises from the conversational context rather than the instruction itself, and we systematically vary the strength of competing evidence by scaling the number of hardcoded turns.

\paragraph{Adversarial context and jailbreaking.}
Jailbreaking attacks construct adversarial contexts that override safety training by exploiting architectural and training vulnerabilities \citep{zou2023universaltransferableadversarialattacks, wei2023jailbrokendoesllmsafety}. Many-shot jailbreaking \citep{anil_many-shot_2024} provides perhaps the closest analog to our setting: it uses hundreds of in-context examples of harmful compliance to induce further compliance. This phenomenon follows predictable power laws where the probability of succumbing to the established pattern increases with the number of contextual turns ($N$) . Other multi-turn strategies, such as Crescendo attacks, achieve similar results through a ``slow boil'' approach \citep{russinovich2025great} starting from initially harmless queries.  We abstract away the core mechanism of escalating in-context evidence of a behavior into a controlled setting with a varied set of instructions of non-harmful parametric control over induction strength ($N$).

\paragraph{In-context exemplars versus instructions.}
A parallel line studies how in-context \emph{exemplars} shape behavior and interact with a model's priors. \citet{minRethinkingRoleDemonstrations2022} show that the input--label format of demonstrations matters more than their correctness, and \citet{weiLargerLanguageModels2023} find that larger models will override semantic priors to fit flipped-label demonstrations---the same induction-over-prior mechanism we isolate. Spurious or irrelevant context degrades adherence \citep{shiLargeLanguageModels2023}, and false demonstrations propagate through the network via late-layer ``false induction heads'' \citep{halawiOverthinkingTruth2024}. Our work differs in placing an \emph{explicit instruction} in conflict with the in-context evidence and scaling that evidence parametrically, rather than studying exemplar quality alone.

\paragraph{Adversarial in-context attacks and instruction hierarchy.}
Beyond many-shot jailbreaking, adversarial or hijacking demonstrations \citep{wangAdversarialDemonstrationAttacks2023, zhouHijackingLargeLanguage2023} and prompt-injection attacks \citep{perezIgnorePreviousPrompt2022} exploit the same porosity to context. A complementary direction trains or evaluates an explicit \emph{instruction hierarchy} that privileges trusted instructions over in-context input \citep{wallaceInstructionHierarchy2024, muCanLLMsFollow2023, zverevCANLLMSSEPARATE2025}---the inverse of the failure we measure, in which in-context evidence overrides a stated instruction.


\paragraph{Self-knowledge and self-prediction.}
Our self-prediction setting connects to a growing body of work on whether language models have privileged access to their own internal states. Prior studies have shown that models finetuned to predict their own behavior outperform cross-model baselines \citep{binderLookingInwardLanguage2024}, though they still lag behind humans in recognizing their own behavioral tendencies \citep{laineMeMyselfAI2024}. More recently findings suggest that models may display emergent introspective abilities, such as distinguishing their own text from artificial prefills \citep{lindseyEmergentIntrospectiveAwareness2026}, as well as articulating implicit behavioral policies that were never explicitly described in training data \citep{betleyTELLMEYOURSELF, plunkettSelfInterpretabilityLLMsCan2025}. We extend these works by testing prospective self-prediction in a zero-shot, dynamic setting, in which models must predict their sensitivity to the instruction-induction conflict.
